% 中文版特有的设置：书名、题注名、定理名、各章引言、页脚。由 main-zh.tex 在 style/book.tex 之后 \input。
% 英文版的对应设置在 style/lang-en.tex，两个文件里的条目一一对应，改了一个，记得看看另一个。

\title{教猫狗学编程（C Language: A Lazy Approach）}
\hypersetup{pdftitle={教猫狗学编程（C Language: A Lazy Approach）}, pdfauthor={whensungoesdown}}

\SetupFloatingEnvironment{listing}{name=代码, listname=代码清单, within=chapter}

% 定理环境，按章编号
\newtheorem{theorem}{定理}[chapter]
\newtheorem{definition}[theorem]{定义}
\newtheorem{lemma}[theorem]{引理}
\renewcommand{\proofname}{证明}

% 页脚：偶数页英文书名，奇数页中文书名，章首页两个都写
\newcommand\lazyfooteven{C Language: A Lazy Approach}
\newcommand\lazyfooterodd{教猫狗学编程}
\newcommand\lazyfooteropener{C Language: A Lazy Approach\quad 教猫狗学编程}

% 每章章首页的引言：\lazyquote{章号}{中文}{原文}{出处}
\lazyquote{1}{学习一门新的程序设计语言的唯一途径，就是用它来编写程序。}
             {The only way to learn a new programming language is by writing programs in it.}
             {Brian W. Kernighan、Dennis M. Ritchie，\\\emph{The C Programming Language}，1978}
\lazyquote{2}{程序是写给人读的，只是顺便让机器执行。}
             {Programs must be written for people to read, and only incidentally for machines to execute.}
             {Harold Abelson、Gerald Jay Sussman，\\\emph{Structure and Interpretation of Computer Programs}，1985}
\lazyquote{3}{人人都知道，调试一个程序比写它难一倍。所以，如果你写程序时已经用尽了聪明才智，又该怎么去调试它呢？}
             {Everyone knows that debugging is twice as hard as writing a program in the first place.
              So if you're as clever as you can be when you write it, how will you ever debug it?}
             {Brian W. Kernighan、P. J. Plauger，\\\emph{The Elements of Programming Style}，第 2 版，1978}
\lazyquote{4}{计算机科学中的所有问题，都可以通过再加一层间接来解决。}
             {All problems in computer science can be solved by another level of indirection.}
             {David J. Wheeler\\（Butler Lampson 在图灵奖演讲中引用）}
\lazyquote{5}{这就是 Unix 哲学：写只做一件事、并且把它做好的程序；写能互相协作的程序。}
             {This is the Unix philosophy: Write programs that do one thing and do it well.
              Write programs to work together.}
             {Doug McIlroy，见 Peter H. Salus，\\\emph{A Quarter Century of Unix}，1994}
\lazyquote{A}{我在做一个（免费的）操作系统（只是业余爱好，不会像 GNU 那样又大又专业），用在 386（486）AT 兼容机上。}
             {I'm doing a (free) operating system (just a hobby, won't be big and professional like gnu)
              for 386(486) AT clones.}
             {Linus Torvalds，新闻组 comp.os.minix，\\1991 年 8 月 25 日}
\lazyquote{B}{这是一次阻止战争的尝试。我希望还不算太晚，希望和平能够——也许是奇迹般地——重新降临。}
             {This is an attempt to stop a war. I hope it is not too late and that somehow,
              magically perhaps, peace will prevail again.}
             {Danny Cohen，\emph{On Holy Wars and a Plea for Peace}，\\1980 年 4 月 1 日}
\lazyquote{C}{任何一个傻瓜都能写出计算机看得懂的代码。好的程序员写的是人看得懂的代码。}
             {Any fool can write code that a computer can understand.
              Good programmers write code that humans can understand.}
             {Martin Fowler，\\\emph{Refactoring}，1999}
